\section{Guide to the main proofs and verification scope}

The proofs that follow use one common chain of ingredients. The signed causal bridge in \cref{obj:prop:signed-causal-bridge} reduces the nonsmooth part of welfare to the average absolute treatment contrast and connects the causal and paired experiments. The private-moment and finite-calibration results in \cref{obj:lem:private-moment-and-dependence,obj:lem:finite-private-polynomial-calibration} establish noninteractive attainment. The moment-duality, adaptive-contraction, and decision lemmas in \cref{obj:lem:cai-low-moment-duality,obj:lem:adaptive-moment-contraction,obj:lem:separated-value-mixtures} supply the sequential converse. The two-cell result in \cref{obj:thm:two-cell-calibration} covers the dimension-free calibration needed across the full resource domain.

The welfare \cref{obj:thm:uniform-private-value-frontiers} combines these ingredients once for both squared-error risk and globally honest connected-interval length. Conditional on the cited Bernstein absolute-approximation limit, the paired \cref{obj:thm:paired-uniform-tv-frontier} transfers the same statistical analysis through the signed experiment, and the full-simplex \cref{obj:thm:full-simplex-tv-converse} then uses inclusion of the paired family and restriction of simplex honesty to that subfamily. The detailed proofs below retain the exact protocol, interval, and resource scopes of those statements.

\subsection{Verification scope}

The formal theorem statements assume uniform strata, fair conditional assignment, consistency, interior arm means, independent participants, and independent public and analyst randomness, as specified in \cref{obj:def:causal-model,obj:ass:iid-people,obj:ass:independent-randomness}. Privacy is pure full-record local privacy; sequential and noninteractive protocols satisfy their respective history conditions in \cref{obj:def:sequential-class,obj:def:noninteractive-class}. Lean verifies the deductions from these hypotheses, including the finite-resource comparisons, protocol transfers, and decision reductions.

The exact asymptotic premise \(2k e_{2k}\to\beta_*\), where \(e_K\) is the best degree-\(K\) uniform approximation error for \(|x|\) on \([-1,1]\), is cited from \citet[Section 3.1]{CaiLow2011Nonsmooth} and \citet{Bernstein1913} rather than formalized. The three theorem-local disclosures identify every main result conditional on this premise. Lean checks how that inverse-degree approximation scale is used and verifies the remaining deductions; it does not certify the cited approximation theorem itself.
